\section{The single-link law of SU(3)}\label{app:su3}

This appendix\footnote{Throughout
this appendix the Lie algebra $\mathfrak{su}(3)$ carries the metric
$\ip XY=-\Tr(XY)$ of the defining representation, the irreducible
representations are labelled by their highest weight $(p,q)$ with respect
to the fundamental weights, of dimension $\frac12(p+1)(q+1)(p+q+2)$, and the
Casimir operator is the Laplace--Beltrami operator of the bi-invariant
metric acting on matrix coefficients; the defining representation is
$(1,0)$ and the adjoint is $(1,1)$.} proves the one-link estimate used in
Theorem~\ref{thm:bpweighted}(ii). Throughout $G=\SU(3)$, the distance is
$d(U,V)=\|U-V\|_{\rm HS}$, which is at most the geodesic distance of the
metric $\ip XY=-\Tr(XY)$ on $\mathfrak{su}(3)$, $K$ is the closed convex hull
of $\SU(3)$ in $\C^{3\times3}$, $A(M)=\frac12(M-M^*)$ is the anti-Hermitian
part of a matrix and $P$ removes the trace, $PX=X-\frac13(\Tr X)I$. With this
metric $\Ric=\frac32$, the Casimir operator acts on the irreducible
representation of highest weight $(p,q)$ by
\begin{equation}\label{eq:casimir}
C(p,q)=\tfrac23(p^2+q^2+pq+3p+3q),
\end{equation}
so that $\lambda_F=C(1,0)=\frac83$ on the fundamental representation and
$C(1,1)=6$ on the adjoint~\cite[Sec.~15.3]{FH91}.

\subsection{Reduction to one link}

For $b\in\C^{3\times3}$ let $\nu_b(dU)\propto e^{\ReTr(Ub)}dU$ on $\SU(3)$. In
the setting of Theorem~\ref{thm:bpweighted}(ii) the conditional law of a free
link $e$ is $\nu_b$ with $b=N\beta\sum_{p\ni e}S_p$, where $S_p$ is the staple
of $e$ in $p$, a product of three matrices each of which is an element of
$\SU(3)$ or the identity, or the adjoint of such a product; so $b\in rK$ with
$r=6N\beta=2\beta_W$. For a matrix $h$ put $\varphi_h(U)=\ReTr(Uh)$ and
\[
L_h(b)=\sup\{\Cov_{\nu_b}(g,\varphi_h):g\text{ $1$-Lipschitz for }d\},\qquad
\Lambda(r)=\sup_{b\in rK}\ \sup_{\|h\|_{\rm HS}=1}L_h(b).
\]
Changing a partner $k$ of $e$ by $d(Q_k,Q'_k)$ changes one staple by a matrix
of Hilbert--Schmidt norm $d(Q_k,Q_k')$, since multiplication by unitaries
preserves that norm, hence $b$ by $h$ with $\|h\|_{\rm HS}\le N\beta\,d(Q_k,
Q'_k)$. Along the segment $b_s=b+sh$, $0\le s\le1$, which stays in $rK$ by
convexity, $\frac d{ds}\nu_{b_s}(g)=\Cov_{\nu_{b_s}}(g,\varphi_h)$, so
by Kantorovich--Rubinstein duality
\begin{equation}\label{eq:su3onelink}
C^T_{ek}\le N\beta\,\Lambda(r)=\frac r6\,\Lambda(r).
\end{equation}
Since $K$ is convex and contains $0$, the mean of $I$, $\omega I$ and
$\omega^2I$ with $\omega=e^{2\pi i/3}$, $r'K\subset rK$ for $r'\le r$ and
$\Lambda$ is nondecreasing. The rest of the appendix bounds $\Lambda(r)$.

\subsection{Gradients and one-forms}

\begin{lemma}\label{lem:su3gradients}
In the left trivialisation, $\nabla\varphi_c(U)=-PA(cU)$; if $V=\varphi_b$
then $\Hess V(X,X)=\ReTr(UX^2b)$; and $\Delta\varphi_c=-\lambda_F\varphi_c$.
\end{lemma}

\begin{proof}
$\frac d{dt}\ReTr(Ue^{tX}c)|_{t=0}=\ReTr(cUX)=-\ip{PA(cU)}X$ for $X\in
\mathfrak{su}(3)$, because $\ReTr(MX)=\ReTr(A(M)X)=-\ip{A(M)}X$ and $X$ is
traceless. On a group with a bi-invariant metric $\nabla_XX=0$ for a
left-invariant field $X$, so the Hessian is the second derivative along
$t\mapsto Ue^{tX}$. The function $\varphi_c$ is a real part of a matrix
coefficient of the fundamental representation, on which the Casimir acts by
$\lambda_F$.
\end{proof}

\begin{lemma}\label{lem:su3rough}
The bottom of the spectrum of the rough Laplacian $\nabla^*\nabla$ on $L^2$
one-forms of $\SU(3)$ is $\mu_0=\frac76$.
\end{lemma}

\begin{proof}
Trivialise one-forms on the left, so that a one-form is a function
$f:\SU(3)\to\mathfrak{su}(3)$. For an orthonormal basis $X_a$ acting as
left-invariant vector fields, $\nabla_{X_a}Y=\frac12[X_a,Y]$ and
$\nabla_{X_a}X_a=0$, so $\nabla^*\nabla=-\sum_a(\rho(X_a)+\frac12\operatorname{ad}
X_a)^2$, with $\rho$ the right regular representation. On the isotypic part
$V_\pi\otimes V_\pi^*\otimes\mathfrak{su}(3)$ of $L^2\otimes\mathfrak{su}(3)$,
given by the Peter--Weyl theorem, put $A_a=\rho_\pi(X_a)$, $B_a=\operatorname{ad}
X_a$ and $T_a=A_a+B_a$, the action on the tensor product. Since $A_a$ and
$B_a$ commute, $A_a+\frac12B_a=\frac12(A_a+T_a)$ and $A_aT_a+T_aA_a=A_a^2+T_a^2
-B_a^2$, so on a component $\sigma\subset\pi\otimes\operatorname{ad}$
\[
\nabla^*\nabla=\tfrac12(C_\pi+C_\sigma)-\tfrac14C_{\rm ad}=\tfrac12(C_\pi+
C_\sigma)-\tfrac32 ,
\]
with $C$ as in~\eqref{eq:casimir}. The components of $(p,q)\otimes(1,1)$
are $(p+1,q+1)$, $(p+2,q-1)$, $(p-1,q+2)$, $(p+1,q-2)$, $(p-2,q+1)$,
$(p-1,q-1)$ and $(p,q)$, the last with multiplicity two, one or zero as both,
one or none of $p,q$ are positive, and only those with nonnegative labels
occur \cite[Sec.~13.2]{FH91}. For $\pi=(0,0)$ the only component is $(1,1)$,
with value $\frac32$. For $\pi=(1,0)$ the components are $(2,1)$, $(0,2)$,
$(1,0)$, with values $\frac{31}6$, $\frac{19}6$, $\frac76$, and $(0,1)$ is
conjugate. For $p+q\ge2$ the smallest values occur at $(p-1,q-1)$, $(p,q)$,
$(p-2,q+1)$ or $(p+1,q-2)$, and each is at least $\frac13(2(p^2+q^2+pq)+3(p+q)
-3)-\frac32\ge\frac32$.
\end{proof}

\subsection{The gap of the one-link measure}

\begin{lemma}\label{lem:su3gap}
For $b\in rK$ the spectral gap of $L=\Delta+\ip{\nabla V}{\nabla\cdot}$,
$V=\varphi_b$, on $L^2(\nu_b)$ is at least
\[
\gamma(r)=\tfrac{39}{64}\mu_0+\tfrac32-\tfrac r2=\tfrac{283}{128}-\tfrac r2 .
\]
\end{lemma}

\begin{proof}
Let $-Lf=\lambda f$ with $\lambda>0$, $\omega=df$ and $\eta=e^{V/2}\omega$;
norms are in $L^2(dU)$. Bochner's formula integrated against $\nu_b$ gives
$\lambda\nu_b|\omega|^2=\nu_b|\nabla\omega|^2+\nu_b((\Ric-\Hess V)(\omega,
\omega))$, and since $e^{V/2}\nabla\omega=\nabla\eta-\frac12dV\otimes\eta$,
\begin{equation}\label{eq:su3bochner}
\lambda\|\eta\|^2=\|\nabla\eta-\tfrac12dV\otimes\eta\|^2+\int(\Ric-\Hess V)
(\eta,\eta).
\end{equation}
Put $\mathcal A=\nabla\eta$ and $\mathcal B=\frac12dV\otimes\eta$. Then
\[
\ip{\mathcal A}{\mathcal B}=\tfrac14\int\ip{\nabla V}{\nabla|\eta|^2}
=\tfrac{\lambda_F}4\int V|\eta|^2 ,
\]
because $-\Delta V=\lambda_FV$ by Lemma~\ref{lem:su3gradients}. Writing
$\|\mathcal A-\mathcal B\|^2=\|\frac58\mathcal A-\mathcal B\|^2+\frac{39}{64}
\|\mathcal A\|^2-\frac34\ip{\mathcal A}{\mathcal B}$ and dropping the first
term,
\[
\|\mathcal A-\mathcal B\|^2\ge\tfrac{39}{64}\mu_0\|\eta\|^2-\tfrac12\int
V|\eta|^2
\]
by Lemma~\ref{lem:su3rough}. With $\Ric=\frac32$, \eqref{eq:su3bochner} gives
$\lambda\|\eta\|^2\ge(\frac{39}{64}\mu_0+\frac32)\|\eta\|^2-\int(\Hess V(\eta,
\eta)+\frac12V|\eta|^2)$, and the claim follows from the pointwise bound
\begin{equation}\label{eq:su3pointwise}
\ReTr(MX^2)+\tfrac12\ReTr(M)\,|X|^2\le\tfrac r2|X|^2,\qquad M=bU\in rK ,
\end{equation}
since $\Hess V(X,X)=\ReTr(MX^2)$ and $V=\ReTr M$. The maximum over $|X|=1$ of
the left side is convex in $M$, so it suffices to take $M=rW$ with $W\in
\SU(3)$ and, after a conjugation, $W=\operatorname{diag}(e^{i\theta_k})$ with
$\sum_k\theta_k=0$; put $c_k=\cos\theta_k$. Write $X=iH$ with $H$ Hermitian,
traceless, $\Tr H^2=1$, with diagonal $d_k$ and entries $H_{jk}$. Then
\[
\ReTr(WX^2)=-\sum_kc_kd_k^2-\sum_{j<k}(c_j+c_k)|H_{jk}|^2 ,
\]
with $\sum_kd_k^2+2\sum_{j<k}|H_{jk}|^2=1$, so the maximum is the larger of
$q_{\rm off}=\max_{j<k}-\frac12(c_j+c_k)$ and $q_{\rm diag}=\max\{-\sum c_kd_k^2:
\sum d_k=0,\ \sum d_k^2=1\}$. If $l$ is the third index, $q_{\rm off}+\frac12
\sum_kc_k=\frac12c_l\le\frac12$. The value $q_{\rm diag}$ is minus the smaller
eigenvalue of $\operatorname{diag}(c)$ compressed to the plane $\sum d_k=0$;
the eigenvalues of the compression solve $\sum_k(c_k-\mu)^{-1}=0$, that is
$3\mu^2-2e_1\mu+e_2=0$ with $e_1,e_2$ the elementary symmetric functions of
the $c_k$, so $q_{\rm diag}=\frac13(\sqrt{e_1^2-3e_2}-e_1)$. The inequality
$q_{\rm diag}+\frac12e_1\le\frac12$ is $\sqrt{e_1^2-3e_2}\le\frac12(3-e_1)$, that
is $e_1^2-4e_2+2e_1-3\le0$. When $\sum\theta_k=0$, $\sum c_k^2=1+2c_1c_2c_3$,
and with this identity the left side equals $-2(1-c_1)(1-c_2)(1-c_3)\le0$.
\end{proof}

The weight $\frac58$ is the one that gives the coefficient $\frac12$ of $V$,
and~\eqref{eq:su3pointwise} is an equality at $W=I$.

\subsection{Haar integrals}

For $h\in\C^{3\times3}$ let $E_2(h)=\frac12((\Tr h)^2-\Tr h^2)$, the sum
of the principal $2\times2$ minors of $h$. All integrals are over $\SU(3)$
with Haar measure; put $t=\Tr U$,
\[
Z_*(r)=\int e^{r\ReTr U}dU,\qquad m_k=\int(\ReTr U)^kdU .
\]

\begin{lemma}\label{lem:su3haar}
\begin{enumerate}[label=\textup{(\roman*)},leftmargin=*]
\item The integral of a monomial of degree $j$ in the entries of $U$ and
$k$ in those of $\bar U$ vanishes unless $j\equiv k$ modulo three;
$\int u_{ij}\bar u_{kl}dU=\frac13\delta_{ik}\delta_{jl}$ and $\int u_{ab}
u_{cd}u_{ef}dU=\frac16\varepsilon_{ace}\varepsilon_{bdf}$.
\item For every $h$ the integrals of $\Tr(hUhU)$ and $(\Tr hU)^2$
vanish, $\int|\Tr hU|^2dU=\frac13\|h\|_{\rm HS}^2$,
\[
\int\Tr(hUhU)\Tr U\,dU=-\tfrac13E_2(h),\qquad
\int(\Tr hU)^2\Tr U\,dU=\tfrac13E_2(h),
\]
and $|E_2(h)|\le\|h\|_{\rm HS}^2$.
\item $\int t^a\bar t^bdU$ vanishes unless $a\equiv b$ modulo three, does
not change when $a$ and $b$ are exchanged, and equals $1$ at $(a,b)=
(0,0)$, $(1,1)$, $(3,0)$, $2$ at $(2,2)$, $3$ at $(4,1)$, $5$ at $(6,0)$,
$6$ at $(3,3)$, $11$ at $(5,2)$, $21$ at $(7,1)$ and $23$ at $(4,4)$.
Hence $m_0=1$, $m_1=0$,
\[
(m_2,\dots,m_8)=\bigl(\tfrac12,\tfrac14,\tfrac34,\tfrac{15}{16},
\tfrac{65}{32},\tfrac{231}{64},\tfrac{973}{128}\bigr),
\]
and $|m_k|\le3^{k-8}m_8$ for $k\ge8$.
\end{enumerate}
\end{lemma}

\begin{proof}
(iii) By Schur--Weyl duality \cite[Thm.~6.3]{FH91}, $t^a=\sum_{\lambda\vdash
a}f^\lambda s_\lambda(U)$ with $f^\lambda$ the number of standard Young
tableaux of shape $\lambda$, and on $\SU(3)$ the Schur polynomial $s_\lambda$
vanishes when $\lambda$ has more than three rows and is otherwise the
irreducible character of highest weight $(\lambda_1-\lambda_3,\lambda_2-
\lambda_3)$. By Schur orthogonality~\cite[Ch.~2]{Bum13}, $\int t^a\bar
t^bdU$ is the sum of $f^\lambda f^\mu$ over the pairs $\lambda\vdash a$,
$\mu\vdash b$ with at most three rows and $\lambda-\mu$ a multiple of
$(1,1,1)$; such pairs exist only when $a\equiv b$ modulo three, and the sum
is symmetric in $a,b$. For $a\ge b$ and $a+b\le8$ the pairs are
$((1),(1))$; $((1,1,1),\varnothing)$; $\lambda=\mu\in\{(2),(1,1)\}$;
$((2,1,1),(1))$; $((2,2,2),\varnothing)$; $\lambda=\mu\in\{(3),(2,1),
(1,1,1)\}$; $((3,1,1),(2))$ and $((2,2,1),(1,1))$; $((3,2,2),(1))$; and
$\lambda=\mu\in\{(4),(3,1),(2,2),(2,1,1)\}$. The hook length formula gives
$f^{(2,1)}=2$, $f^{(3,1)}=f^{(2,1,1)}=3$, $f^{(2,2)}=2$, $f^{(2,2,2)}=5$,
$f^{(3,1,1)}=6$, $f^{(2,2,1)}=5$ and $f^{(3,2,2)}=21$, hence the values, the
last one being $1+9+4+9$. Since $(\ReTr U)^k=2^{-k}\sum_a\binom kat^a\bar
t^{k-a}$,
\begin{gather*}
m_2=\tfrac{2}{4},\quad m_3=\tfrac28,\quad m_4=\tfrac{6\cdot2}{16},\quad
m_5=\tfrac{2\cdot5\cdot3}{32},\\
m_6=\tfrac{2\cdot5+20\cdot6}{64},\quad m_7=\tfrac{2\cdot21\cdot11}{128},
\quad m_8=\tfrac{2\cdot8\cdot21+70\cdot23}{256}.
\end{gather*}
The eigenvalues of $U$ have modulus one, so $|\ReTr U|\le3$ and
$|\ReTr U|^k\le3^{k-8}(\ReTr U)^8$ for $k\ge8$.

(i) The map $U\mapsto\omega U$, $\omega=e^{2\pi i/3}$, preserves Haar
measure and multiplies the monomial by $\omega^{j-k}$. For a unitary
representation $\pi$ of $\SU(3)$, $\int\pi(U)dU$ is the orthogonal
projection onto the invariant vectors, whose dimension is the mean of the
character of $\pi$. On $\C^3\otimes\overline{\C^3}$ the invariant vectors
are the multiples of the identity, which gives the second formula. On
$(\C^3)^{\otimes3}$ the character is $t^3$, whose mean is $1$ by (iii), and
$\varepsilon$ is invariant because $\det U=1$; the projection onto the line
of $\varepsilon$, whose squared norm is $6$, gives the third formula.

(ii) The first two integrals have degree two in $U$ and none in $\bar U$,
and $\int|\Tr hU|^2dU=\frac13\sum|h_{ji}|^2$ by (i). By (i),
$\int h_{pa}u_{aq}h_{qc}u_{cp}u_{ee}dU=\frac16h_{pa}h_{qc}\varepsilon_{ace}
\varepsilon_{qpe}$, summed over repeated indices, and $\sum_e
\varepsilon_{ace}\varepsilon_{qpe}=\delta_{aq}\delta_{cp}-\delta_{ap}
\delta_{cq}$, which gives $\frac16(\Tr h^2-(\Tr h)^2)$; the second integral
is the same with $\varepsilon_{pqe}$ in place of $\varepsilon_{qpe}$, which
changes the sign. Finally $E_2(h)$ is the trace of the second exterior power
of $h$, whose singular values are the products $\sigma_i\sigma_j$, $i<j$, of
those of $h$; so $|E_2(h)|\le\sum_{i<j}\sigma_i\sigma_j=\frac12((\sum_i
\sigma_i)^2-\sum_i\sigma_i^2)\le\sum_i\sigma_i^2$.
\end{proof}

\subsection{The modulus}

\begin{lemma}\label{lem:su3modulus}
Let $r\ge0$ with $\gamma(r)>0$, put
\[
\epsilon(r)=1-\frac{5\,e^{-r/5}}{9\,Z_*(r)},
\]
and let $b\in rK$ and $\|h\|_{\rm HS}=1$. Then
\[
\nu_b|\nabla\varphi_h|^2\le\epsilon(r),\quad\Var_{\nu_b}(\varphi_h)\le
\frac{\epsilon(r)}{\gamma(r)},\quad L_h(b)\le\frac{\sqrt{\epsilon(r)}}
{\gamma(r)},
\]
so that $\Lambda(r)\le\sqrt{\epsilon(r)}/\gamma(r)$ and, by
\eqref{eq:su3onelink}, $C^T_{ek}\le r\sqrt{\epsilon(r)}/(6\gamma(r))$. The
function $\epsilon$ is nondecreasing on $[0,\infty)$.
\end{lemma}

\begin{proof}
By Lemma~\ref{lem:su3gradients}, $\nabla\varphi_h(U)=-PA(M)$ with $M=hU$, and
$\|M\|_{\rm HS}=1$. The anti-Hermitian part has $|A(M)|^2=\frac12(\|M\|^2-
\ReTr M^2)$, and $P$ removes from it the multiple $\frac13\Tr A(M)=\frac i3
\operatorname{Im}\Tr M$ of the identity, whose squared norm is
$\frac13(\operatorname{Im}\Tr M)^2$. Hence
\begin{gather*}
|\nabla\varphi_h|^2=1-\tfrac12F_h,\\
F_h(U)=1+\ReTr(hUhU)+\tfrac23(\operatorname{Im}\Tr hU)^2 ,
\end{gather*}
and $F_h\ge0$ because $|\ReTr M^2|\le\|M\|^2$. For complex numbers
$X,Y$, $\operatorname{Re}X\operatorname{Re}Y=\frac12\operatorname{Re}(XY)+
\frac12\operatorname{Re}(X\bar Y)$ and $(\operatorname{Im}X)^2=\frac12|X|^2-
\frac12\operatorname{Re}X^2$. With these identities, the vanishing of the
integrals of degree two in $U$ and one in $\bar U$, or the converse, and
Lemma~\ref{lem:su3haar}(ii),
\[
\int F_h\,dU=1+\tfrac23\cdot\tfrac16=\tfrac{10}9,\quad\int F_h\ReTr U\,dU=
-\tfrac29\operatorname{Re}E_2(h),
\]
the second as the sum of $-\frac16\operatorname{Re}E_2(h)$ and $\frac23
\cdot(-\frac1{12})\operatorname{Re}E_2(h)$. For $W\in\SU(3)$ the
substitution $U\mapsto W^{-1}U$ and Lemma~\ref{lem:su3haar}(ii) give
\[
\int F_h\ReTr(WU)\,dU=\int F_{hW^{-1}}\ReTr U\,dU\ge-\tfrac29 .
\]

Write $b=r\sum_i\lambda_iW_i$ as a convex combination of elements of
$\SU(3)$ and $V=\ReTr(Ub)$. Under the probability measure $\pi=\frac9{10}
F_h\,dU$ the mean of $V$ is at least $-\frac9{10}\cdot\frac29r=-\frac r5$,
so Jensen's inequality gives $\int F_he^VdU=\frac{10}9\pi(e^V)\ge\frac{10}9
e^{-r/5}$. The function $c\mapsto\log\int e^{\ReTr(Uc)}dU$ is convex by
H\"older's inequality and equals $\log Z_*(r)$ at every $c=rW$, $W\in
\SU(3)$, by the invariance of Haar measure; hence $\int e^VdU\le Z_*(r)$
and $\nu_b(F_h)\ge\frac{10}9e^{-r/5}/Z_*(r)$, which is the first bound.

The spectral gap of $\nu_b$ is at least $\gamma(r)$ by Lemma~\ref{lem:su3gap}.
The Poincar\'e inequality gives $\Var_{\nu_b}(\varphi_h)\le\nu_b|\nabla
\varphi_h|^2/\gamma(r)$, the second bound, and, for $g$ $1$-Lipschitz for
$d$, hence for the geodesic distance, $\Var_{\nu_b}(g)\le1/\gamma(r)$; so
$\Cov_{\nu_b}(g,\varphi_h)\le\Var_{\nu_b}(g)^{1/2}\Var_{\nu_b}(\varphi_h)^{1/2}
\le\sqrt{\epsilon(r)}/\gamma(r)$. Finally $\log Z_*$ is convex in $r$ with
derivative $m_1=0$ at $r=0$, so $Z_*$ and $\epsilon$ are nondecreasing on
$[0,\infty)$.
\end{proof}

The losses in this bound are Jensen's inequality, the Poincar\'e inequality
for $\varphi_h$, and the gap $\gamma(r)$ itself: at $b=0$ the variance of
$\varphi_h$ is $\frac16$ by Lemma~\ref{lem:su3haar}(ii) and the gap is
$\lambda_F=\frac83$, while the lemma uses $\epsilon(0)=\frac49$, which is
exact, and $\gamma(0)=\frac{283}{128}$, so at $b=0$ only the gap is lost.
